\newpage
\pagestyle{plain}

\section{\centering РЕЗУЛЬТАТЫ}
\subsection{Оптимальные значения коэффициентов и значения ошибки}
\subsubsection{Детерминированный сеточный поиск}
На этапе грубого поиска по регулярной сетке параметров с шагом $\Delta c$ алгоритм вычисляет значения целевого функционала только для 72 конфигураций, порождённых условием монотонности \eqref{eq:monotony}. Остальные $120$ точек, не соответствующие условию, отбрасываются ещё до вызова интегратора.

После завершения этого этапа алгоритм переходит к трёхстадийному измельчению шага вокруг найденного локального минимума. На первом этапе уточнения шаг равен $0,04$, на втором — $0,02$,  на третьем - $0,01$. \autoref{tab:opt_results_grid_search} приводит значения ошибки для каждой из задач на всех этапах сеточного поиска.

\begin{table}
	\centering
	\caption{Результаты оптимизации сеточным поиском}
	\label{tab:opt_results_grid_search}
	\vspace{5mm}
	\small
	
	\renewcommand{\arraystretch}{1.5}
	
	\begin{tabular}{|c|c|c|c|c|}
		\hline
		\textbf{Задача} &
		\textbf{Грубый поиск} &
		\textbf{Изм. 1} &
		\textbf{Изм. 2} &
		\textbf{Изм. 3}
		\\ \hline
		
		Лин. осциллятор &
		$6.531 \cdot 10^{-7}$ &
		$2.389 \cdot 10^{-7}$ &
		$2.389 \cdot 10^{-7}$ &
		$1.224 \cdot 10^{-7}$
		\\ \hline
		
		Нелин. уравнение Шрёдингера &
		$2.763 \cdot 10^{-4}$ &
		$2.763 \cdot 10^{-4}$ &
		$1.435 \cdot 10^{-4}$ &
		$9.265 \cdot 10^{-5}$
		\\ \hline
		
		Задача двух тел &
		$1.203 \cdot 10^{0}$ &
		$1.203 \cdot 10^{0}$ &
		$1.174 \cdot 10^{0}$ &
		$1.143 \cdot 10^{0}$
		\\ \hline
		
		Задача пяти внешних тел &
		$1.554 \cdot 10^{-9}$ &
		$1.430 \cdot 10^{-9}$ &
		$2.241 \cdot 10^{-10}$ &
		$4.369 \cdot 10^{-11}$
		\\ \hline
		
	\end{tabular}
	
	\renewcommand{\arraystretch}{1}
\end{table}

\subsubsection{Модифицированный метод роя частиц}
Зафиксируем значения шага $h$ и временные отрезки задачи Коши для каждой из четырёх тестовых задач:
\begin{itemize}
	\item Лин. осциллятор: $h = 0.1$, $t_0 = 0$, $T= 10$;
	\item Нелин. уравнение Шрёдингера: $h = 0.05$, $t_0 = 0$, $T = 2$;
	\item Задача двух тел: $h = 0,05$, $t_0 = 0$, $T = 15$;
	\item Задача пяти внешних тел: $h = 0.2$, $t_0 = 0$, $T = 20$.
\end{itemize}

\autoref{tab:opt_results} приводит результаты решений всех четырёх тестовых задач при размере популяции в 30 частиц и ограничением в 30 итераций и позволяет сравнить ошибку оптимизированного метода с ошибкой классического метода, под которым понимается 8-стадийный непараметрический метод Рунге--Кутты шестого порядка. Коэффициенты его таблицы Бутчера задаются вектором параметров $\theta_{\text{cl}} = [0.1, 0.25, 0.4, 0.7]$, который не учитывает осциллирующий или нелинейный характер правых частей уравнений ОДУ.  Можно видеть, что прирост точности различается от задачи к задаче из-за различий в рельефах целевого функционала \eqref{eq:functional}, частоты колебаний и размерности исследуемой системы ОДУ.

\begin{table}
	\centering
	\caption{Результаты оптимизации методом роя частиц}
	\label{tab:opt_results}
	\vspace{5mm} 
	\small
	\renewcommand{\arraystretch}{1.5}
	\begin{tabular}{|c|c|c|c|}
		\hline
		\multicolumn{1}{|c|}{\textbf{Задача}} &
		\textbf{Классич. метод} &
		\textbf{Оптим. метод.} &
		\multicolumn{1}{|c|}{\textbf{Прирост точности}} \\ \hline
		
		Лин. осциллятор &
		$2.800 \cdot 10^{-5}$ &
		$4.092 \cdot 10^{-9}$ &
		$6842.4$ \\ \hline 
		
		Нелин. уравнение Шрёдингера &
		$6.836 \cdot 10^{-4}$ &
		$3.817 \cdot 10^{-5}$ &
		$17.9$ \\ \hline
		
		Задача двух тел &
		$1.637 \cdot 10^{0}$ &
		$2.245 \cdot 10^{-1}$ &
		$7.3$ \\ \hline
		
		Задача пяти внешних тел &
		$1.969 \cdot 10^{-8}$ &
		$3.545 \cdot 10^{-11}$ &
		$555.6$ \\ \hline
	\end{tabular}
	\renewcommand{\arraystretch}{1}
\end{table}

Приведём также векторы оптимальных значений $\theta^*$, полученные методом роя частиц для каждой из задач:
\begin{itemize}
	\item Лин. осциллятор: $[0.07819, 0.35000, 0.47245, 0.98000]$;
	\item Нелин. уравнение Шрёдингера: $[0.09010, 0.23948, 0.42150, 0.87405]$;
	\item Задача двух тел: $[0.13642, 0.21989, 0.38210, 0.81445]$;
	\item Задача пяти внешних тел: $[0.10806, 0.22604, 0.41023, 0.97626]$.
\end{itemize}

Обратим внимание на парадокс результатов задачи двух тел: ошибка решения задачи оптимизированным методом составляет $2.245 \times 10^{-1}$, в то время как ошибка решения более сложной задачи пяти внешних тел имеет порядок $10^{-11}$. Такое поведение ошибки объясняется различиями в характерах начальных условий и геометрии траекторий движения. В задаче пяти внешних тел моделируется движение планет-гигантов, которые совершают почти круговое движение вокруг Солнца. Напротив, в постановке задачи двух тел задана низкая начальная скорость, что приводит к сильно вытянутой эллиптической орбите. При прохождении перицентра расчетная точка резко приближается к притягивающему центру, где ньютоновские гравитационные силы демонстрируют взрывообразный рост пропорционально $1/r^2$. Тем не менее, разработанный оптимизатор смог снизить погрешность решения в $7.3$ раза относительно классической схемы, что подтверждает целесообразность использования алгоритма в жестких расчетных режимах.

Данная гипотеза подтверждается увеличением начальной скорости до единицы, которое переводит систему в режим круговой орбиты постоянного радиуса. Таким образом, правая часть системы ОДУ представляет собой гладкие колебания, а итоговая погрешность метода упала до $6.467 \cdot 10^{-9}$.

\subsubsection{Сравнение полученных результатов с указанными в первоисточниках}
Сравнение полученных с помощью нашей реализации результаты на примере уравнения Шрёдингера демонстрирует расхождение с результатами, полученными в работах \cite{Anastassi2023, Anastassi2025}, что иллюстрирует влияние схемы оптимизации на точку сходимости в многоэкстремальном пространстве параметров.

Так, в работе \cite{Anastassi2023} получен вектор $[0,007; 0,208; 0,441; 0,915]$, а в более поздней работе \cite{Anastassi2025} параметры сместились, дав вектор $[0,009; 0,216; 0.544; 0,980]$. Расходимость с выведенным нами вектором $[0,091; 0,239; 0;421; 0.874]$ вызвана тремя причинами: заменой аналитического решения системы из 29 условий порядка численным, предотвращением вырождения матриц и усилением метода роя частиц с помощью локальной полировки.

\subsection{Визуализация решения}

